\section{Обзор существующих решений}
\subsection{Анализаторы сетевого трафика}
Современные системы анализа сетевого трафика можно классифицировать по глубине инспектирования, архитектуре обработки и модели расширения функциональности. В данном подразделе рассматриваются три наиболее распространённых открытых решения: Ntopng, Suricata и Zeek, с акцентом на их подходы к генерации и обработке событий.

\subsubsection{Ntopng}
Ntopng представляет собой систему мониторинга сетевого трафика, ориентированную на визуализацию потоков, анализ пропускной способности и выявление аномалий на основе статистических моделей. Архитектура Ntopng базируется на потоковом анализе (flow-based), где пакеты агрегируются в записи NetFlow/IPFIX с последующим хранением в Redis или базе данных TimescaleDB. Генерация событий в Ntopng реализована преимущественно через пороговые механизмы и интеграцию с внешними системами оповещения. Расширяемость обеспечивается через плагины на языке Lua и REST API, однако внутренняя событийная шина не выделена в отдельный компонент, что ограничивает гибкость маршрутизации событий и требует перезапуска процессов при глубокой модификации логики обработки. Система эффективна для задач мониторинга производительности, но не предназначена для детального анализа протоколов прикладного уровня в реальном времени \cite{ntopng2023doc}.

\subsubsection{Suricata}
Suricata является многопоточным IDS/IPS-движком, поддерживающим сигнатурный анализ, обнаружение аномалий и глубокую инспекцию пакетов. Архитектура Suricata построена на основе конвейера обработки, где каждый пакет проходит через стадии дефрагментации, реконструкции потоков и применения правил. События генерируются при совпадении сигнатур или нарушении правил безопасности и выводятся в форматах JSON, Eve или syslog. Расширение функциональности возможно через добавление новых правил в формате YAML/Lua и подключение внешних модулей, однако базовый механизм событийного движка жёстко привязан к ядру обработки. Динамическое добавление новых типов событий требует компиляции или перезапуска, а обработка событий осуществляется синхронно в рамках потока обработки пакетов, что создаёт узкие места при высокой нагрузке. Suricata демонстрирует высокую пропускную способность, но её событийная модель не оптимизирована для гибкой аналитической кастомизации без вмешательства в исходный код \cite{suricata2023arch}.

\subsubsection{Zeek}
Zeek (ранее Bro) позиционируется как система мониторинга сетевой безопасности, ориентированная на протокольный анализ и генерацию структурированных логов. В отличие от сигнатурных систем, Zeek использует событийно-ориентированную модель, где каждый этап обработки пакета (установление соединения, запрос/ответ протокола, завершение сессии) генерирует соответствующее событие. Обработка событий осуществляется через встроенный скриптовый язык ZeekScript, позволяющий описывать логику анализа без модификации ядра. Однако интерпретация скриптов вносит значительные накладные расходы, а управление состоянием событий ограничено однопоточной или слабо масштабируемой моделью. Для внесения изменений в схему событий или добавления новых обработчиков часто требуется перезагрузка политик или перезапуск процесса, что снижает доступность в режимах непрерывного мониторинга. Zeek остаётся эталоном в области протокольной диссекции, но её событийная подсистема не удовлетворяет требованиям к высокой пропускной способности и динамической расширяемости в современных распределённых средах \cite{zeek2022design}.

\subsubsection{Итоговый анализ анализаторов}
Сравнительный анализ показывает, что существующие системы либо фокусируются на потоковой статистике (Ntopng), либо на сигнатурном обнаружении (Suricata), либо на скриптовой событийной обработке (Zeek). Ни одно из решений не предоставляет выделенной, высокопроизводительной и типобезопасной библиотеки событий, предназначенной исключительно для встраивания в сторонние DPI-конвейеры. Ограничения включают отсутствие динамического реестра схем событий, синхронную обработку с блокирующими вызовами, необходимость перезапуска для обновления логики и слабую изоляцию пользовательских обработчиков. Эти недостатки формируют технологический пробел, который заполняется разработкой независимой расширяемой системы генерации событий.

\subsection{Библиотеки систем событий}
Для реализации событийной логики в программных системах часто применяются готовые библиотеки, обеспечивающие механизмы pub/sub, обработку сигналов и асинхронную диспетчеризацию. В данном подразделе рассматриваются три наиболее релевантные библиотеки в контексте высокопроизводительных систем.

\subsubsection{Boost.Signal2}
Boost.Signal2 представляет собой библиотеку, реализующую механизм сигналов и слотов для языка C++. Она обеспечивает типобезопасную связь между источниками событий и обработчиками, поддерживает множественные соединения, приоритеты и группировку слушателей. Библиотека использует виртуальные вызовы и динамическое выделение памяти для управления соединениями, что обеспечивает гибкость, но вносит значительные накладные расходы на этапе диспетчеризации. В контексте обработки сетевого трафика, где требуется генерация миллионов событий в секунду с предсказуемой задержкой, использование Boost.Signal2 приводит к фрагментации памяти, увеличению времени отклика из-за виртуализации вызовов и отсутствию оптимизаций под кэш-локальность. Библиотека подходит для приложений общего назначения, но не удовлетворяет требованиям к low-latency системам реального времени \cite{boost2023signal}.

\subsubsection{Eventpp}
Eventpp — современная C++-библиотека, предоставляющая легковесный событийный движок с поддержкой компиляционной типобезопасности, лямбда-выражений и асинхронной обработки. В отличие от Boost, Eventpp минимизирует использование виртуальных функций и динамических аллокаций, применяя шаблоны и интрузивные списки для хранения подписок. Библиотека демонстрирует высокую производительность в микробенчмарках и подходит для встраиваемых систем и игровых движков. Однако Eventpp не включает механизмов маршрутизации событий между потоками, управления пулами обработки, изоляции сбоев или версионирования схем. В контексте сетевого анализа отсутствие встроенной поддержки lock-free очередей и NUMA-оптимизаций ограничивает её применимость в распределённых DPI-архитектурах. Eventpp служит хорошей базой, но требует значительной доработки для промышленного использования в высоконагруженных конвейерах \cite{eventpp2023doc}.

\subsubsection{Libuv}
Libuv представляет собой кроссплатформенную библиотеку асинхронного ввода-вывода, лежащую в основе Node.js. Она реализует событийный цикл (event loop), предоставляет механизмы обработки файловых дескрипторов, таймеров, потоков и IPC. Libuv оптимизирована для высокой пропускной способности сетевых операций и поддерживает масштабируемую модель обработки через пулы потоков. Однако её абстракция ориентирована на асинхронный I/O и управление ресурсами операционной системы, а не на типобезопасную генерацию и маршрутизацию структурных событий внутри приложения. Интеграция Libuv в DPI-конвейер потребовала бы дополнительного слоя абстракции для типизации событий, управления схемами и обеспечения детерминированной диспетчеризации, что снижает эффективность и увеличивает сложность кодовой базы. Libuv решает задачи асинхронности, но не покрывает требования к событийной модели анализа трафика \cite{libuv2023doc}.

\subsubsection{Итоговый анализ библиотек событий}
Существующие библиотеки либо обеспечивают высокую гибкость за счёт производительности (Boost.Signal2), либо оптимизированы под специфические домены без поддержки сетевых паттернов (Eventpp, Libuv). Отсутствие комплексного решения, сочетающего компиляционную типобезопасность, асинхронную lock-free диспетчеризацию, масштабируемость на многоядерных архитектурах и расширяемую схему событий, подтверждает необходимость разработки специализированной библиотеки. Разрабатываемая система должна устранить выявленные недостатки, предоставив унифицированный API, оптимизированный для встраивания в DPI-конвейеры и обеспечения предсказуемого поведения под нагрузкой.

\newpage